\documentclass[12pt, a4paper, oneside]{ctexart}
\usepackage{amsmath, amsthm, amssymb, bm, xcolor, tikz, framed, graphicx, hyperref, mathrsfs}
\usepackage{bookmark}
\usepackage{pgfplots}
\pgfplotsset{compat=1.18}

\newcommand{\bra}[1]{\langle #1 |}
\newcommand{\ket}[1]{| #1 \rangle}
\newcommand{\braket}[2]{\langle #1 | #2 \rangle}

\title{Quantum Mechanics Homework (Week 14)}
\author{Tang Yi \\ [学号已脱敏]}
\usepackage[a4paper, left=2cm, right=2cm, top=2cm, bottom=2cm]{geometry}
\linespread{1.3} % 行距更紧凑
\definecolor{shadecolor}{RGB}{241, 241, 255}

\newcounter{problemname}
\newenvironment{problem}
  {\begin{shaded}\stepcounter{problemname}\par\noindent\textbf{题目 \theproblemname. }}
  {\end{shaded}\par}
\newenvironment{solution}
  {\par\noindent\textbf{题目 \theproblemname 的解答. }}
  {\par}
\newenvironment{note}
  {\par\noindent\textbf{题目 \theproblemname 的注记. }}
  {\par}

\begin{document}

\maketitle


\begin{problem}
  % w15-1
考虑三维希尔伯特空间的哈密顿量：
$$H = V_0 \begin{pmatrix} (1-\epsilon) & 0 & 0 \\ 0 & 1 & \epsilon \\ 0 & \epsilon & 2 \end{pmatrix}$$
其中 $V_0$ 为常数，$\epsilon$ 为小量。

1) 求出无微扰 ($\epsilon = 0$) 时的本征态和本征值；

2) 严格求解 $H$ 的本征值，结果展开为 $\epsilon$ 的幂级数，保留到二阶；

3) 求 $H^0$ 的非简并本征值在二级修正下的结果，并与精确结果比较；
\end{problem}

\begin{solution}
  \[
H = V_0
\begin{pmatrix}
1-\epsilon & 0 & 0 \\
0 & 1 & \epsilon \\
0 & \epsilon & 2
\end{pmatrix},
\qquad \epsilon \ll 1
\]

(1) 无微扰情形 $\epsilon=0$。

此时哈密顿量为
\[
H^{(0)} = V_0
\begin{pmatrix}
1 & 0 & 0 \\
0 & 1 & 0 \\
0 & 0 & 2
\end{pmatrix}.
\]

其本征值与对应本征态为
\[
E_1^{(0)} = V_0, \quad |1\rangle = (1,0,0)^T,
\]
\[
E_2^{(0)} = V_0, \quad |2\rangle = (0,1,0)^T,
\]
\[
E_3^{(0)} = 2V_0, \quad |3\rangle = (0,0,1)^T.
\]

其中 $E_1^{(0)}$ 与 $E_2^{(0)}$ 构成二重简并能级。

\bigskip

(2) 严格求解 $H$ 的本征值并展开到 $\epsilon^2$ 阶。

由于态 $|1\rangle$ 与其余态完全解耦，其本征值可直接得到
\[
E_1 = V_0(1-\epsilon).
\]

其余两个能级由 $2\times2$ 子矩阵
\[
V_0
\begin{pmatrix}
1 & \epsilon \\
\epsilon & 2
\end{pmatrix}
\]
决定。

特征方程为
\[
\lambda^2 - 3\lambda + (2-\epsilon^2)=0,
\]
其解为
\[
\lambda_{\pm}=\frac{3\pm\sqrt{1+4\epsilon^2}}{2}.
\]

对应能量
\[
E_{\pm}=V_0\lambda_{\pm}.
\]

对平方根作小量展开
\[
\sqrt{1+4\epsilon^2}=1+2\epsilon^2+O(\epsilon^4),
\]
得到
\[
E_2 = V_0(1-\epsilon^2)+O(\epsilon^4),
\]
\[
E_3 = V_0(2+\epsilon^2)+O(\epsilon^4).
\]

\bigskip

(3) 非简并本征值的二级微扰修正及比较。

将哈密顿量分解为
\[
H = H^{(0)} + H',
\qquad
H' = V_0
\begin{pmatrix}
-\epsilon & 0 & 0 \\
0 & 0 & \epsilon \\
0 & \epsilon & 0
\end{pmatrix}.
\]

考虑非简并态 $|3\rangle$。

一阶能量修正为
\[
E_3^{(1)}=\langle 3|H'|3\rangle=0.
\]

二阶能量修正为
\[
E_3^{(2)}=\sum_{n\neq3}\frac{|\langle n|H'|3\rangle|^2}
{E_3^{(0)}-E_n^{(0)}}.
\]

只有 $n=2$ 项有贡献，
\[
\langle 2|H'|3\rangle=V_0\epsilon,
\qquad
E_3^{(0)}-E_2^{(0)}=V_0,
\]
因此
\[
E_3^{(2)}=V_0\epsilon^2.
\]

最终得到
\[
E_3 = 2V_0 + V_0\epsilon^2,
\]
与精确本征值的展开结果
\[
E_3 = V_0(2+\epsilon^2)
\]
完全一致。

\end{solution}

\begin{problem}
  % w15-2
设非简谐振子的哈密顿量表示为 $H = H_0 + H'$，其中
$$H_0 = -\frac{\hbar^2}{2\mu} \frac{\mathrm{d}^2}{\mathrm{d}x^2} + \frac{1}{2}\mu\omega^2 x^2$$
$$H' = \beta x^3 \quad (\beta \text{ 为常数})$$
用微扰法求其能量本征值（准到二级近似）和本征函数（准到一级近似）。
\end{problem}

\begin{solution}
设
\[
H = H_0 + H', \qquad H' = \beta x^3 ,
\]
其中
\[
H_0 = -\frac{\hbar^2}{2\mu}\frac{\mathrm{d}^2}{\mathrm{d}x^2}
+ \frac{1}{2}\mu\omega^2 x^2 .
\]

无微扰谐振子的本征值与本征态为
\[
E_n^{(0)} = \hbar\omega\left(n+\frac{1}{2}\right),
\qquad
|n\rangle ,
\quad n=0,1,2,\dots
\]

引入升降算符
\[
x = \sqrt{\frac{\hbar}{2\mu\omega}}\,(a+a^\dagger),
\]
因此
\[
x^3=\left(\frac{\hbar}{2\mu\omega}\right)^{3/2}
(a+a^\dagger)^3 .
\]


一阶修正为
\[
E_n^{(1)}=\langle n|H'|n\rangle
=\beta\langle n|x^3|n\rangle .
\]

由于 $x^3$ 为奇函数，而 $|n\rangle$ 具有确定宇称，
\[
\langle n|x^3|n\rangle=0,
\]
因此
\[
E_n^{(1)}=0.
\]


二阶修正公式为
\[
E_n^{(2)}=\sum_{m\neq n}
\frac{|\langle m|H'|n\rangle|^2}
{E_n^{(0)}-E_m^{(0)}}.
\]

展开
\[
(a+a^\dagger)^3
= a^3 + 3a^2a^\dagger + 3aa^{\dagger 2} + a^{\dagger 3},
\]
可知 $x^3$ 只连接
\[
m = n\pm1,\quad n\pm3 .
\]

计算矩阵元并代入能量差，整理后得到
\[
E_n^{(2)}
= -\frac{\beta^2\hbar^2}{8\mu^3\omega^4}
\left(30n^2+30n+11\right).
\]

\end{solution}

\begin{problem}
  % w15-3
接 w15-1，4) 求 $H^0$ 的简并本征值在一级修正下的结果，并与精确结果比较。
\end{problem}

\begin{solution}
  (4) $H^{(0)}$ 的简并本征值在一级修正下的结果及与精确结果的比较。

无微扰哈密顿量
\[
H^{(0)} = V_0
\begin{pmatrix}
1 & 0 & 0 \\
0 & 1 & 0 \\
0 & 0 & 2
\end{pmatrix}
\]
在能量
\[
E^{(0)} = V_0
\]
处存在二重简并，对应简并子空间由
\[
|1\rangle = (1,0,0)^T, \qquad
|2\rangle = (0,1,0)^T
\]
张成。

微扰项为
\[
H' = V_0
\begin{pmatrix}
-\epsilon & 0 & 0 \\
0 & 0 & \epsilon \\
0 & \epsilon & 0
\end{pmatrix}.
\]

将 $H'$ 限制在简并子空间 $\{|1\rangle,|2\rangle\}$ 上，得到有效微扰矩阵
\[
H'_{\text{deg}}
=
V_0
\begin{pmatrix}
\langle 1|H'|1\rangle & \langle 1|H'|2\rangle \\
\langle 2|H'|1\rangle & \langle 2|H'|2\rangle
\end{pmatrix}
=
V_0
\begin{pmatrix}
-\epsilon & 0 \\
0 & 0
\end{pmatrix}.
\]

对该矩阵对角化，得到一级能量修正
\[
\Delta E_1^{(1)} = -V_0\epsilon,
\qquad
\Delta E_2^{(1)} = 0.
\]

因此，简并能级在一级微扰下劈裂为
\[
E_1 = V_0(1-\epsilon),
\qquad
E_2 = V_0.
\]

对应的一级本征态仍为
\[
|1\rangle = (1,0,0)^T,
\qquad
|2\rangle = (0,1,0)^T.
\]


与精确结果的比较：

在 w15-1 第 (2) 问中，精确对角化得到能量展开为
\[
E_1^{\text{exact}} = V_0(1-\epsilon),
\]
\[
E_2^{\text{exact}} = V_0(1-\epsilon^2) = V_0 + O(\epsilon^2).
\]

可见，在一级近似下
\[
E_1^{(1)} = V_0(1-\epsilon)
\]
与精确结果完全一致，而
\[
E_2^{(1)} = V_0
\]
与精确结果在 $O(\epsilon)$ 阶一致，其差异出现在二级修正。

因此，简并微扰论在一级近似下正确给出了简并能级的劈裂结构。

\end{solution}
\begin{problem}
  % w15-4
粒子被限制在圆环内自由转动，其哈密顿量为：
$$H = \frac{\hat{l}_z^2}{2I}$$
若在圆环的其中一段加上弱常势场：$V(\varphi) = V_0, -\beta < \varphi < \beta$，求基态和第一激发态能量的一级修正。
\end{problem}

\begin{solution}
  粒子在圆环上自由转动，其无微扰哈密顿量为
\[
H_0=\frac{\hat l_z^2}{2I},
\qquad
\hat l_z=-i\hbar\frac{\partial}{\partial\varphi}.
\]

无微扰本征方程为
\[
H_0\psi_m(\varphi)=E_m^{(0)}\psi_m(\varphi),
\]
其本征值与本征函数为
\[
E_m^{(0)}=\frac{\hbar^2 m^2}{2I},
\qquad
\psi_m(\varphi)=\frac{1}{\sqrt{2\pi}}e^{im\varphi},
\quad m=0,\pm1,\pm2,\dots
\]

其中基态为 $m=0$，第一激发态为 $m=\pm1$，具有二重简并。

加入微扰势
\[
V(\varphi)=
\begin{cases}
V_0, & -\beta<\varphi<\beta,\\
0, & \text{其它区域}.
\end{cases}
\]


(1) 基态能量的一级修正。

一级能量修正为
\[
E_0^{(1)}=\langle 0|V|0\rangle
=\int_{-\pi}^{\pi}\psi_0^*(\varphi)\,V(\varphi)\,\psi_0(\varphi)\,\mathrm d\varphi.
\]

代入基态波函数
\[
\psi_0(\varphi)=\frac{1}{\sqrt{2\pi}},
\]
得到
\[
E_0^{(1)}
=\frac{1}{2\pi}\int_{-\beta}^{\beta}V_0\,\mathrm d\varphi
=\frac{V_0}{2\pi}(2\beta)
=\frac{V_0\beta}{\pi}.
\]


(2) 第一激发态能量的一级修正。

第一激发态 $m=\pm1$ 构成简并子空间，
\[
\psi_{\pm1}(\varphi)=\frac{1}{\sqrt{2\pi}}e^{\pm i\varphi}.
\]

计算微扰矩阵元
\[
V_{mn}=\langle m|V|n\rangle
=\frac{V_0}{2\pi}\int_{-\beta}^{\beta}e^{i(n-m)\varphi}\,\mathrm d\varphi.
\]

在 $\{ |1\rangle, |-1\rangle \}$ 基底下，有
\[
V_{11}=V_{-1,-1}=\frac{V_0\beta}{\pi},
\]
\[
V_{1,-1}=V_{-1,1}
=\frac{V_0}{2\pi}\int_{-\beta}^{\beta}e^{-2i\varphi}\,\mathrm d\varphi
=\frac{V_0}{2\pi}\frac{\sin(2\beta)}{1}.
\]

因此简并子空间内的有效微扰矩阵为
\[
V_{\text{deg}}=
\begin{pmatrix}
\dfrac{V_0\beta}{\pi} & \dfrac{V_0}{2\pi}\sin(2\beta) \\
\dfrac{V_0}{2\pi}\sin(2\beta) & \dfrac{V_0\beta}{\pi}
\end{pmatrix}.
\]

对角化该矩阵，得到一级能量修正
\[
E_{\pm1}^{(1)}
=\frac{V_0\beta}{\pi}
\pm \frac{V_0}{2\pi}\sin(2\beta).
\]


基态能量在一级近似下为
\[
E_0 = 0 + \frac{V_0\beta}{\pi}.
\]

第一激发态的简并在一级微扰下被解除，其能量为
\[
E_{1,\pm}
=\frac{\hbar^2}{2I}
+\frac{V_0\beta}{\pi}
\pm \frac{V_0}{2\pi}\sin(2\beta).
\]

\end{solution}
\begin{problem}
  % w15-5
同样是圆环内的自由转子，带电情况下其电偶极矩 $\mathbf{D} = -q\mathbf{r}$，若加上与圆环平行，强度为 $\mathcal{E}$ 的匀强电场，分别计算：

1) 电场很弱时，一级修正下的本征能量和本征函数；

2) 电场很强时，转子只能局限在一个很窄的角度内振动，此时的基态能量。
\end{problem}

\begin{solution}
  粒子带电量 $q$，在半径为 $R$ 的圆环上运动，其电偶极矩为
\[
\mathbf D = -q\mathbf r.
\]

设匀强电场 $\boldsymbol{\mathcal E}$ 与圆环平行，可取其方向为
\[
\boldsymbol{\mathcal E} = \mathcal E\,\hat x .
\]

在极角 $\varphi$ 表示下，
\[
\mathbf r = R(\cos\varphi,\sin\varphi),
\]
因此微扰势能为
\[
H' = -\mathbf D\cdot\boldsymbol{\mathcal E}
= q\mathcal E R\cos\varphi .
\]

无微扰哈密顿量及其本征态为
\[
H_0=\frac{\hat l_z^2}{2I},
\qquad
E_m^{(0)}=\frac{\hbar^2 m^2}{2I},
\qquad
\psi_m(\varphi)=\frac{1}{\sqrt{2\pi}}e^{im\varphi},
\quad m=0,\pm1,\pm2,\dots
\]


(1) 电场很弱时的一级微扰结果。

一级能量修正为
\[
E_m^{(1)}=\langle m|H'|m\rangle
=q\mathcal E R
\int_{-\pi}^{\pi}\frac{\mathrm d\varphi}{2\pi}
e^{-im\varphi}\cos\varphi\,e^{im\varphi}.
\]

由于
\[
\int_{-\pi}^{\pi}\cos\varphi\,\mathrm d\varphi = 0,
\]
可得
\[
E_m^{(1)}=0,
\qquad \forall m.
\]

因此在一级近似下，所有能级的能量均不发生移动。

一级近似下的本征函数修正为
\[
|m^{(1)}\rangle
=\sum_{n\neq m}
\frac{\langle n|H'|m\rangle}
{E_m^{(0)}-E_n^{(0)}}\,|n\rangle.
\]

将
\[
\cos\varphi=\frac{1}{2}(e^{i\varphi}+e^{-i\varphi})
\]
代入，可得选择定则
\[
n = m\pm1.
\]

矩阵元为
\[
\langle m\pm1|H'|m\rangle
=\frac{q\mathcal E R}{2}.
\]

因此本征函数在一级近似下为
\[
|m\rangle_{\text{pert}}
=|m\rangle
+\frac{q\mathcal E R}{2}
\left(
\frac{|m+1\rangle}{E_m^{(0)}-E_{m+1}^{(0)}}
+\frac{|m-1\rangle}{E_m^{(0)}-E_{m-1}^{(0)}}
\right).
\]


(2) 电场很强时的近似基态能量。

强电场极限下，势能项
\[
V(\varphi)=q\mathcal E R\cos\varphi
\]
在 $\varphi=0$ 处取得最小值。

在 $\varphi\simeq0$ 附近展开
\[
\cos\varphi\simeq1-\frac{\varphi^2}{2},
\]
势能近似为
\[
V(\varphi)\simeq
q\mathcal E R
-\frac{1}{2}q\mathcal E R\,\varphi^2.
\]

忽略常数项后，有效哈密顿量为
\[
H\simeq
-\frac{\hbar^2}{2I}\frac{\mathrm d^2}{\mathrm d\varphi^2}
+\frac{1}{2}q\mathcal E R\,\varphi^2.
\]

这是一个一维谐振子，其有效角频率满足
\[
\frac{1}{2}I\omega_{\text{eff}}^2
=\frac{1}{2}q\mathcal E R,
\]
即
\[
\omega_{\text{eff}}=\sqrt{\frac{q\mathcal E R}{I}}.
\]

因此基态能量近似为
\[
E_0\simeq
-q\mathcal E R
+\frac{1}{2}\hbar\sqrt{\frac{q\mathcal E R}{I}}.
\]

其中第一项为势能最低点，第二项为零点能。

\end{solution}
\end{document}